\documentclass[
reprint,
superscriptaddress,
frontmatterverbose,
amsmath,amssymb,
pra,
]{revtex4-2}

\usepackage{amsfonts}
\usepackage{amsmath}
\usepackage{amssymb}
\usepackage{amsthm}
\usepackage[normalem]{ulem}
\usepackage{graphicx}
\usepackage{dcolumn}
\usepackage{bm}
\usepackage{listings}
\usepackage{xcolor}
\usepackage{csquotes}
\usepackage{hyperref}
\usepackage{array}
\usepackage{booktabs}
\usepackage{multirow}
\usepackage{braket}

\definecolor{codegreen}{rgb}{0,0.6,0}
\definecolor{codegray}{rgb}{0.5,0.5,0.5}
\definecolor{codepurple}{rgb}{0.58,0,0.82}
\definecolor{backcolour}{rgb}{0.95,0.95,0.92}

\lstdefinestyle{mystyle}{
    backgroundcolor=\color{backcolour},
    commentstyle=\color{codegreen},
    keywordstyle=\color{magenta},
    numberstyle=\tiny\color{codegray},
    stringstyle=\color{codepurple},
    basicstyle=\ttfamily\footnotesize,
    breakatwhitespace=false,
    breaklines=true,
    captionpos=b,
    keepspaces=true,
    numbers=left,
    numbersep=5pt,
    showspaces=false,
    showstringspaces=false,
    showtabs=false,
    tabsize=2
}
\lstset{style=mystyle}

\hypersetup{
    pdfborder={0 0 0},
    pdfnewwindow=true,
    colorlinks=true,
    hypertexnames=false,
    linkcolor=blue,
    citecolor=blue,
    filecolor=blue,
    urlcolor=blue,
}

\newcommand{\ii}{{\mathrm{i}}}
\newcommand{\bk}{{\bm{k}}}
\newcommand{\bR}{{\bm{R}}}
\newcommand{\tr}{{\text{tr}}}
\renewcommand\boldsymbol{\bm}

\begin{document}

\title{VBCSR: High-Performance Distributed Infrastructure for Materials Simulation}

\author{Zhanghao Zhouyin}
\affiliation{Department of Physics, McGill University, Montreal, Quebec, Canada H3A2T8}

\author{Hong Guo}
\affiliation{Department of Physics, McGill University, Montreal, Quebec, Canada H3A2T8}

\date{\today}

\begin{abstract}
Modern materials simulation is increasingly built from locality: localized atomic orbitals, elements, descriptors, learned Hamiltonian blocks etc.  These local objects are natural for chemistry, physics, and machine learning, but in many fields, simulation using a global operator build by its local parts is essential. Therefore, a highly-scalable, distributed operator infrastructure that have simple dependencies, high-performance and convinent enough for python dominant usage is desired. VBCSR is a Python-first distributed operator infrastructure designed for this need.  It represents atom-centered and graph-structured operators as sparse matrices with scalar, fixed-block, or variable-block structure. For maximize performance, VBCSR automatically selects highly tuned CSR, BSR, or true variable-block kernels behind a common user interface.  The package combines distributed data structure supported for general development, as well as the domain specific functionals that is right out-of-the-box. This paper mainly introduce the design concept, the architecture and main functionalities. We hope this package can help people benefit from the simple usage so to build their scientific workflow easier, while still benefiting the high-performance design to wrote simple but productive code with minimal effort to assist AI-enabled and first-principles and materials simulation.
\end{abstract}

\maketitle

\section{Introduction}

Locality principles are widely adopted in modern materials simulation.  In localized-basis electronic structure, a Hamiltonian or density matrix is built from atom-centered orbital blocks.  In tight-binding and machine-learned Hamiltonian models, the natural input and output are interactions on a local neighbor graph.  In finite-element, finite-volume, and coupled multi-physics methods, space is discretized into local cells, elements, or nodes that may carry different numbers of unknowns.  These examples come from different communities, but they lead to the same computational object when scaling up: a distributed infrastructure that represents local couplings between spatially localized degrees of freedom.

\begin{figure}[t]
    \centering
    \includegraphics[width=0.92\linewidth]{fig/logo.png}
    \caption{The VBCSR package logo.}
    \label{fig:logo}
\end{figure}

The locality principle have largely motivated the AI-enabled scientific computing.  Machine Learning interatomic potentials, material descriptors, Hamiltonian blocks and surrogate properties mostly prefer local representations since they often embeded the dominant interactions, designing in this principle often lead to transferbility and better accuracy.  Many of the practical simulation, however, must do more than predict local pieces.  It must assemble those pieces into a global operator, distribute the operator across many processes, apply it to vectors state or batches of states, evaluate matrix functions, and often requires superhigh performance beyond prototype Python workflows.  Available packages often sacrifice the simplicity in exchange of the performance, despite highly mature, would less ideal for daily research usage.

VBCSR is designed around one rule: expose a simple scientific object as Python native, and hide the complicated backend choices needed for performance.  The user works with graphs, atoms, local blocks, vectors, multivectors, and operator algebra just like operating any other Python object, while benefitting from the highly tunned kernels for efficiency and scalability. Therefore, the user can easily implement prototype algorithm with VBCSR, or replace their original code's data structure with VBCSR with performance boost. For developers, VBCSR are designed to be a lightweight, dependency-minimal, and composable package that can be embedded in larger workflows.

The VBCSR stores and applies is centred around the structure of operators. Specifically, the operator is stored in the VBCSR as sparse matrices through CSR, BSR or truly variable-block CSR kernels to match the diverse data layout. The variable-block case is essential for atom-centered simulation. For example, in density functional theory and tight-binding simulations, different species naturally carry different numbers of orbitals; a carbon atom, a hydrogen atom, and a transition-metal atom should be considered as the element in the spatial domain that defines the interconnectivity patterns, while the local degree of freedom should vary. Similar structures is also common in finite element simulations.  More generally, a node $i$ may carry $b_i$ local degrees of freedom, and a coupling to node $j$ is a dense block of shape $b_i \times b_j$.  we design specific storage, computational kernels and  the distribution framework to naturally benefit from this structure. Therefore, VBCSR can becomes useful to a broad scientific audience where the performance is crucial.

Existing high-performance sparse and electronic-structure libraries address important parts of this landscape, including RESCU \cite{michaud2016rescu}, DBCSR \cite{sivkov2019dbcsr}, CP2K \cite{kuhne2020cp2k}, NTPoly \cite{dawson2018massively}, BML \cite{bock2018basic}, CheSS \cite{mohr2017efficient}, and specialized dense micro-kernel approaches such as LIBXSMM \cite{heinecke2016libxsmm}.  VBCSR benefit from the mature design of algorithms of those, while concentrate on the niche: providing a Python-first distributed infrastructure layer that can be embedded in scientific workflows in material simulation  with multiple high-performance backend.  Therefore, a scalar graph Laplacian, a uniform-block elasticity operator, and a species-dependent localized-basis Hamiltonian can therefore share one user-facing abstraction.

The contribution of this paper is consequently both numerical and infrastructural.  We present VBCSR as a lightweight distributed operator package for materials simulation, with:
\begin{enumerate}
    \item graph, vector, multivector, and matrix abstractions that preserve
    spatial-local and graph-centered degrees of freedom;
    \item automatic CSR, BSR, and VBCSR backend switching under one Python API;
    \item shape-aware variable-block storage and batched kernels for sparse apply,
    sparse-dense apply, thresholded sparse multiplication and graph-local matrix functions; and
    \item atomistic and periodic-image containers that connect localized-basis
    materials data to distributed operator algebra.
\end{enumerate}

\section{Lightweight Python Operator Infrastructure}

\subsection{Distributed operators as central object}

The central abstraction in VBCSR is a distributed operator defined on a graph. The local degree of freedom in each node of the graph defines the block size of the operator, and the edges of the graph defines the sparsity pattern. The spatial distribution is therefore a graph level partition using either geometric inspired method, or graph theory based method. Therefore, the structure and operator themselves can be defined distributedly by their local blocks, and automatically assembled to form the global representation. Worthnoticing is, the graph and partition is automatically conducted, and the logic is hidden from the usage, therefore, the user can focus on the scientific problem and constructing the operators by adding local blocks, without worrying anything about the data geometry and distribution.

Once the structure of the partitioned graph are defined, the DOFs of the nodes are used to define the rows and columns of the operator.  Each node $i$ carries a block size $b_i$, and an operator block between nodes $i$ and $j$ is stored as a dense matrix $A_{ij}\in \mathbb{C}^{b_i \times b_j}$ or $A_{ij}\in \mathbb{R}^{b_i \times b_j}$.  Applying the operator has the form
\begin{equation}
    y_i = \sum_{j \in \mathcal{N}(i)} A_{ij} x_j,
    \label{eq:block_apply}
\end{equation}
where $x_j$ and $y_i$ are local vector or blocks of vectors, and $\mathcal{N}(i)$ is the
neighbor set of node $i$. Therefore, we can constructed the oprator class in the scientific object like the Hamiltonian in atomic simulation, the differential operator in PDE solving, and use the operator operation to solve it by iterative methods. 

The repo also have other structures, they are all design as object that act by, auxilliate, or extend the operator class. For example, the disvector and multivector are defined as the distributed state with assist the operator, while the atomic data and image container are defined as the high-level scientific object to assist the atomic system construction, which is backed by the distribution class of the operator, and the real space and reciprocal space operator representation and conversion.  The design is therefore a layered infrastructure, where the user can use the high-level scientific object to construct the operator, and use the operator to solve the problem, while the backend is fully hidden so no complexity is exposed.

\subsection{High Performance}

High performance in VBCSR comes naturally from the data model of VBCSR.  Once the graph and block sizes are known, the package chooses an internal storage format that matches the operator.  Scalar operators use a CSR path, uniformly blocked operators use a BSR path, and genuinely heterogeneous operators use the VBCSR path.  The user still sees one matrix object and the same operations, such as construction, assembly, multiplication, transposition, filtering, and matrix functions.

This automatic choice matters because different scientific operators expose
different kinds of regularity.  A scalar graph descriptor should be stored and
applied as a scalar sparse matrix.  A model with the same number of unknowns on
every node can use fixed-size dense blocks.  A multi-species localized-basis
Hamiltonian may instead contain many block shapes, because different atoms
carry different numbers of orbitals.  VBCSR keeps these cases under one
interface, but lowers them to different kernels behind the scenes.

The storage layout is designed to keep the numerical work close to the memory
layout.  The CSR and BSR backends use paged sparse storage and can call highly
optimized sparse kernels from vendor libraries when they are available.  The 
variable-block backend is accelerated through a shape resolved structure. We notice that despite the shape of blocks might varies, the total kind of shapes are often finite. Therefore, we groups blocks with the same shape and applies batched dense kernels to those groups when they applied.  In practical terms, the package tries to give the computer long, regular pieces of work despite the scientific operator is irregular.

Parallel performance follows the same locality principle.  The distributed
graph assigns each process a set of owned nodes and records the neighboring
remote nodes that are needed to apply the local block rows.  During an
operator application, only these neighboring vector blocks have to be
exchanged.  The partition can be produced from geometric locality, from a graph
partition, or from user-provided distributed data, so the computational layout
can follow the spatial structure of the physical system. Practically, a partition can be optimized to minimize the neighbour required to achieve maximum performance.  For the scientific
application, this allows the developer fully disentagle with the high-performance computing aspects of the problem. You only need to form the local blocks, and the package will handle the distribution and communication automatically.

\subsection{Python-native workflow}

The Python interface is designed so that a distributed operator behaves like a
normal scientific Python object.  The matrix class follows the SciPy
\texttt{LinearOperator} model and supports familiar matrix syntax, including
the \texttt{@} operator for applying a matrix to a vector, a batch of vectors,
or another sparse operator.  This makes VBCSR usable inside ordinary Python
workflows without requiring users to write separate code for the distributed
case.

Vectors and multivectors are also exposed in a NumPy-like way.  Their locally
owned entries can be viewed or filled with NumPy arrays, so prototype data,
test cases, and small validation examples can move naturally between VBCSR and
the wider scientific Python ecosystem.  On small problems, conversion helpers
allow comparison with SciPy sparse matrices; on large problems, the same
logical operations are executed by the distributed C++ backend.

The key design choice is that advanced operator functionality is embedded in
the object itself.  Scaling, shifting, transposition, conjugation, sparse
matrix multiplication, block filtering, submatrix extraction, multivector
application, and matrix-function approximations are called as ordinary matrix
operations.  The accelerated kernels, communication schedules, and backend
selection remain hidden behind this interface.  For the user, the workflow is
therefore close to standard Python linear algebra, while the implementation
can still use the specialized storage and kernels needed for large materials
simulations.

\section{Architecture and Algorithms}

\begin{figure*}[t]
    \centering
    \includegraphics[width=0.98\linewidth]{Architecture.png}
    \caption{VBCSR uses one Python-facing matrix interface and dispatches to
    CSR, BSR, or variable-block backends according to the block-size structure
    recorded in the distributed graph.  The same graph ownership, vector
    layouts, communication routines, and numerical libraries support all
    backends.}
    \label{fig:architecture}
\end{figure*}

\subsection{Architecture overview}

Figure~\ref{fig:architecture} summarizes the VBCSR software stack.  The input
structure is a distributed block graph: graph edges define the sparse pattern,
and node block sizes define the local degrees of freedom.  A single Python
matrix class exposes assembly, sparse application, transpose, sparse matrix
multiplication, filtering, and matrix-function utilities through pybind11-backed
native kernels.

The central architectural layer is backend abstraction.  After graph
construction, the matrix facade selects one storage family from the global block
structure: scalar graphs use CSR, uniform non-scalar graphs use BSR, and
heterogeneous graphs use VBCSR.  This dispatch separates public matrix semantics
from physical storage.  The same API is preserved across formats, while each
backend owns the layout, page metadata, cached apply plans, and optional vendor
descriptors required by its kernel family.

Kernel dispatch follows the selected backend.  CSR uses scalar sparse kernels,
BSR uses fixed-size dense-block kernels, and VBCSR uses shape-grouped
variable-block kernels.  All backends share the distributed graph, vector and
multivector layouts, ghost exchange routines, MPI communication layer,
thread-level parallelism, and BLAS/LAPACK dense operations.

\subsection{Distributed graphs, vectors, and ownership}

The distributed graph defines row ownership, local indexing, block dimensions,
and communication metadata.  Each rank owns a subset of global block rows; all
remote column blocks required by those rows are stored as ghosts.  Local indices
are ordered as owned blocks followed by ghost blocks, and adjacency arrays use
this local numbering.  After ghost synchronization, sparse application reduces
to local operation of Eq.~\eqref{eq:block_apply}.

All communication and layout management is handled by the graph.  Therefore, the graph also owns the block sizes.  Their prefix sums define scalar offsets for
matrix blocks, vector storage, and MPI message sizes.

All other object, \texttt{DistVector} and \texttt{DistMultiVector}, uses the graph schedule to exchange owned block data into the ghost region; adjoint operations reverse the schedule and accumulate ghost contributions into owned entries.

\subsection{Variable-block storage}

The VBCSR backend targets graphs with heterogeneous block sizes.  It avoids both
largest-block padding and scalar expansion by separating graph order from value
order.  The graph preserves the node-level sparsity pattern, while numerical
blocks are grouped by shape.

For an edge $(i,j)$, the value block has shape $b_i \times b_j$.  Backend
construction registers each distinct shape, counts its local blocks, and stores
all blocks of that shape in a packed value stream.  A graph-block pointer list maps
each adjacency entry to its location in the shape store, preserving sparsity while giving kernels shape-contiguous workload.

Shape storage is paged, so does the storage for CSR and BSR in the package.  Each shape class owns bounded pages containing
contiguous same-shape blocks, avoiding large monolithic allocations and keeping
index ranges bounded for large matrices.

The shaped pages gives direct benefit for batched kernel execution. Continuous buffers can be hand over copyfree and use highly optimized batched GEMM kernels for supported matrix operation. This gives us superior efficiency, and almost naive compatibilities for GPU acceleration.

\subsection{Sparse apply and sparse-dense apply}

Sparse application is evaluated on the distributed block graph.  After
synchronizing ghost input blocks, each rank applies the owned block rows as
\begin{equation}
    Y_i = \sum_{j \in \mathcal{N}(i)} A_{ij} X_j ,
    \label{eq:block_apply_local}
\end{equation}
where $X_j$ and $Y_i$ may be vectors or dense blocks containing multiple
right-hand sides.  The adjoint apply uses the transposed graph action
\begin{equation}
    Y_j \mathrel{+}= A_{ij}^{\dagger} X_i,\qquad j\in\mathcal{N}(i),
    \label{eq:block_apply_adjoint}
\end{equation}
with ghost-column contributions reduced back to their owning ranks.  Vector and
multivector paths therefore share the communication schedule; the local kernel
changes from GEMV-like to GEMM-like work.

Kernel selection follows the active backend.  CSR can use native sparse kernels
or vendor sparse handles when MKL or AOCL-Sparse is available.  BSR uses vendor
BSR kernels when enabled and otherwise dispatches common small uniform block
sizes to specialized dense-block kernels.  VBCSR traverses same-shape pages and
executes them either as individual dense block products or as strided batched
GEMV/GEMM through the linked BLAS.  Small packed blocks use hand-written
compile-time kernels, including AVX2/FMA paths for double precision; larger or
non-packed cases fall back to CBLAS.  The current implementation does not
include a LIBXSMM-specific path.  OpenMP parallelism is applied over sparse rows
or VBCSR batches, with per-thread accumulation in the variable-block case.

\subsection{Thresholded sparse-sparse multiplication}

Sparse matrix-matrix multiplication appears in density-matrix purification,
polynomial expansions, graph closure operations, and approximate inverse
construction.  VBCSR follows the block-sparse thresholding principle used in
DBCSR \cite{sivkov2019dbcsr}: products are generated at the block level,
\begin{equation}
    C_{ij} = \sum_k A_{ik} B_{kj},
    \label{eq:spgemm}
\end{equation}
and norm estimates are used to avoid forming blocks whose contribution is below
a prescribed tolerance.  This reflects a common principle in linear-scaling
electronic structure: locality and decay make many formally nonzero matrix
entries scientifically negligible
\cite{prodan2005nearsightedness,benzi2013decay,benzi2007decay}.

The multiplication is separated into symbolic and numeric phases.  In the
symbolic phase, each output row accumulates candidate columns using estimates
$\|A_{ik}\|\|B_{kj}\|$; columns whose accumulated estimate exceeds the
threshold are retained in the result graph.  For distributed input, only
metadata for remote rows of $B$ is exchanged during this phase, and the
corresponding remote block values are fetched only after the surviving
structure is known.  The numeric phase evaluates the retained block products,
applies a row-scaled product skip, and performs a final block-norm filter on
the assembled result.

The main difference from a fixed-block multiplication is the backend-adaptive
representation.  CSR, BSR, and VBCSR share the same symbolic filtering design,
while their numeric phases use kernels matched to the active storage format.
The VBCSR path additionally groups products by the row, inner, and column block
dimensions, so heterogeneous blocks are executed as same-shape dense-product
batches rather than padded to a single global block size.  The output is a new
distributed graph and matrix in the same backend family as the inputs.

\subsection{Graph-local matrix functions}

Matrix functions can be approximated by graph-local methods when the relevant
operator is local or rapidly decaying on a sparse graph.  Following this
principle, as in graph-based purification methods \cite{niklasson2016graph},
VBCSR evaluates $f(A)$ from restricted neighborhoods rather than from a global
dense matrix.  For a node $i$, the package extracts an $\ell$-hop neighborhood
$\mathcal{N}_{\ell}(i)$, evaluates a dense matrix function on the induced local
submatrix, and writes back the central block row:
\begin{equation}
    f(A)_{i,:} \approx
    \left[
    f\left(A_{\mathcal{N}_{\ell}(i),\mathcal{N}_{\ell}(i)}\right)
    \right]_{i,:}.
    \label{eq:local_function}
\end{equation}
The approximation is controlled by the graph cutoff radius $\ell$: increasing
$\ell$ enlarges the local subproblem and reduces the truncation error, while a
fixed cutoff keeps the amount of work per graph node bounded.  The algorithm is
therefore linear-scaling for bounded-degree graphs and fixed $\ell$.  In the
distributed setting, each rank imports the halo blocks needed to form its local
neighborhoods, performs the dense evaluations independently, and inserts the
resulting central rows into the distributed output matrix.

\subsection{Atomistic data and periodic quantum operators}

VBCSR includes a native atomistic data layer so that sparse block structure can
be derived directly from the physical system.  \texttt{AtomicData} stores
atomic numbers, positions, cell vectors, periodic boundary flags, atom types,
and species-dependent orbital counts.  Structures may be constructed from
Python/ASE objects, embedded file readers, explicit graph arrays, or already
distributed atom partitions.  From these inputs the package builds the
distributed neighbor graph, records the lattice shift associated with each
periodic edge, assigns atom-centered block sizes, and exchanges ghost atom
attributes needed by local operator assembly.  The automatic construction can
partition the structure internally, while the distributed constructor preserves
a partition supplied by an external workflow.

Periodic quantum operators are represented by \texttt{ImageContainer}, which
stores a family of block-sparse matrices indexed by lattice translation
$\mathbf{R}$.  Each image has its own distributed graph, but all images share
the same atom ownership and orbital block-size model inherited from
\texttt{AtomicData}.  Operator blocks can be inserted individually or in
batches, assembled across MPI ranks, combined image-wise, and redistributed
between compatible atom partitions.  This design keeps real-space operator
construction local while preserving the information required for reciprocal
space.

Sampling at a wave vector $\mathbf{k}$ forms a complex block-sparse operator on
the base atom graph,
\begin{equation}
    O(\mathbf{k}) =
    \sum_{\mathbf{R}} e^{-2\pi\ii\,\mathbf{k}\cdot\mathbf{R}}
    O(\mathbf{R}),
    \label{eq:k_sampling}
\end{equation}
with an optional basis-position phase convention that replaces
$\mathbf{R}$ by $\mathbf{R}+\boldsymbol{\tau}_j-\boldsymbol{\tau}_i$ for block
$(i,j)$.  Thus the same real-space quantum-operator container supports
efficient construction, redistribution, and $R$-space to $k$-space
transformation without discarding the variable-block atom-centered structure.

\section{Benchmark Experiments}

The benchmark section is organized around two questions.  The first is kernel
efficiency: for a fixed sparse operator, how much performance is gained or lost
by using the CSR, BSR, or VBCSR backend rather than a scalar sparse baseline?
The second is distributed cost: how the same design behaves when graph
construction, communication, memory use, and MPI scaling are included.  This
separation keeps single-node kernel comparisons distinct from the broader
infrastructure claims of the package.

\subsection{Kernel efficiency}

The efficiency benchmark compares three operation classes supported by VBCSR:
sparse matrix-vector multiplication (SpMV), sparse matrix-dense matrix
multiplication (SpMM), and sparse general matrix-matrix multiplication
(SpGEMM).  Each operation is measured in three structural domains.  The CSR
domain uses scalar block size and tests whether VBCSR's abstraction introduces
overhead when no block structure is present.  The BSR domain uses a uniform
non-scalar block size and tests the fixed-block path against scalar expansion.
The VBCSR domain uses heterogeneous block sizes and tests the package's main
target case, where padding to a single block size is undesirable.

For each domain, the same mathematical operator will be benchmarked with three
implementations: SciPy sparse matrices, an MKL sparse backend, and the selected
VBCSR backend.  In the CSR domain, all methods operate on scalar CSR
representations.  In the BSR domain, VBCSR and MKL can use fixed-block storage
when available, while SciPy CSR provides the scalar-expanded reference.  In the
VBCSR domain, VBCSR preserves the original heterogeneous block layout; SciPy
and MKL are evaluated on equivalent scalar CSR matrices because their standard
sparse interfaces do not represent variable block sizes directly.  The primary
reported quantity is time per operation, with speedup normalized against the
best available scalar CSR baseline.  Secondary quantities include matrix
storage, number of scalar nonzeros after expansion, number of block nonzeros,
and, for SpGEMM, the output fill ratio.

SpMV measures the overhead and benefit of block traversal for a single
right-hand side.  SpMM measures reuse of the same sparse graph across multiple
dense right-hand sides, where block kernels and batched execution should become
more important.  SpGEMM measures both symbolic graph construction and numeric
block-product accumulation; it is the most sensitive test of whether the block
structure is preserved through a complete sparse algebra operation.  The final
paper should report these results in one combined figure, using rows for the
three operations and columns for the CSR, BSR, and VBCSR domains.

\begin{figure*}[t]
    \centering
    \fbox{\begin{minipage}[c][0.32\textheight][c]{0.92\textwidth}
    \centering
    Planned combined efficiency benchmark.  Rows: SpMV, SpMM, and SpGEMM.
    Columns: CSR, BSR, and VBCSR structural domains.  Each panel reports
    operation time for SciPy sparse, MKL sparse, and VBCSR, with speedup
    annotations relative to the best scalar CSR baseline.
    \end{minipage}}
    \caption{Planned summary figure for kernel-efficiency benchmarks.  The
    figure will replace the current separate preliminary panels after the
    benchmark suite is rerun under a fixed software and hardware environment.}
    \label{fig:benchmark_efficiency}
\end{figure*}

\subsection{Distributed scaling and reproducibility}

The second benchmark group evaluates the costs that are not visible in a
single-kernel timing.  These experiments should measure strong and weak scaling
for distributed SpMV, SpMM, and SpGEMM; graph construction and matrix assembly
time; ghost exchange volume; memory footprint; and the cost of converting
between atomistic operators and sparse matrix backends.  For thresholded
SpGEMM, the benchmark should additionally report the tolerance, skipped block
products, output block count, and final fill ratio, so that performance can be
read together with the approximation and sparsity behavior.

All benchmark results should be generated from scripts that record the package
revision, compiler, MPI implementation, BLAS and sparse backend, CPU model,
thread count, process count, matrix-generation parameters, and random seed.
The repository tests for public Python symbols, distributed graph and vector
behavior, sparse operations, atomistic utilities, and benchmark drivers provide
the correctness layer for this section.  The final manuscript should therefore
report both timing results and the reproducibility metadata needed to rerun each
figure.

\section{Discussion}

VBCSR should be understood as infrastructure for AI-enabled scientific
computing, not as an AI model by itself.  Its role is to make the output of
local models usable in scalable simulation.  A learned Hamiltonian, a
data-driven preconditioner, or a graph-based surrogate operator still needs a
distributed representation, communication schedule, matrix functions, and
validation path.  VBCSR provides that layer while keeping the user close to the
scientific objects: atoms, neighborhoods, orbital blocks, periodic quantum
operators, and matrix functions.

This positioning also prevents the work from reading as an isolated kernel
study.  The numerical kernels are essential, but they serve a larger scientific
story.  The value comes from making several audiences meet in one software
object.  A materials scientist sees localized orbitals and periodic operators.
An AI-for-science researcher sees a route from local predictions to global
simulation.  A numerical linear algebra researcher sees sparse matrix
functions, thresholded multiplication, and block-aware formats.  A high
performance computing researcher sees MPI ownership, communication reuse, and
backend specialization.  The central narrative remains simple: local scientific
representations need distributed operator infrastructure.

Several limitations guide the benchmark design.  First, the efficiency study
must separate CSR, BSR, and VBCSR structural domains rather than reporting one
mixed comparison.  Second, distributed measurements should include larger
systems, strong scaling, weak scaling, memory use, communication volume, and
threshold-accuracy tradeoffs.  Third, future benchmark sets should include
operators produced by external electronic-structure or learned-model pipelines,
so that synthetic matrix results can be connected to application-scale input
distributions.  These additions will strengthen the infrastructure argument
without changing the core manuscript structure.

\section{Conclusion}

VBCSR provides a Python-first distributed operator layer for materials
simulation and related scientific computing workflows.  Its main idea is
simple: scientific users should be able to work with local blocks, atoms,
graphs, periodic quantum operators, and matrix functions, while the package handles
whether the underlying sparse kernel is CSR, BSR, or true variable-block
storage.  This automatic format selection is not only a performance
convenience; it is what allows one workflow to contain scalar, fixed-block, and
variable-block operators without fragmenting the code.

The current implementation combines distributed graphs, vectors, multivectors,
adaptive sparse matrices, thresholded sparse-sparse multiplication, graph-local
matrix functions, atomistic containers, and periodic quantum-operator sampling.
The benchmark plan isolates the corresponding performance questions through
SpMV, SpMM, and SpGEMM comparisons across CSR, BSR, and VBCSR domains, followed
by distributed scaling and reproducibility measurements.  With these benchmark
results and application-scale inputs, VBCSR can serve as a visible piece of
infrastructure for the next generation of AI-enabled materials simulation.

\section*{Code and data availability}

The scripts for generating the current benchmark data and the VBCSR library are
openly available at \url{https://github.com/DeepElectron/vbcsr}.  At
publication, the exact code revision, benchmark inputs, and plotting scripts
used for the figures should be archived with the manuscript.

\appendix

\bibliography{main.bib}
\end{document}
